% !TEX root = ../tex/1-1.tex
% 学習の定式化と順方向KL
\begin{frame}{1-1-1 Training of DGM}
\begin{block}{学習}
モデル族$\{p_{\phi}\}$の中で不一致$\mathcal{D}$を最小化するパラメータを求める\\
\centering $\phi^{*}\in \argmin_{\phi}\mathcal{D}(p_{data},p_{\phi})$
\end{block}
\begin{itemize}
\item $\boldsymbol{p}_{data}$は、未知であるためi.i.d.サンプルから値を効率的に\\
推定することによって実用的なDを選ぶ\\
\item モデルに十分な表現力があれば$\boldsymbol{p}_{data}$と$\boldsymbol{p}_{\phi}$の近似ができる\\
\end{itemize}

\begin{block}{順方向KL divergenceと最尤推定}
    MLE(Maximum Likelihood Estimation): 最尤推定\\
    ⇒訓練データから最も尤もらしいパラメータ$\phi^{*}$を推定する方法\\
    \small ※モチベーション：モデルの確率密度$p_{\phi}(x)$を評価できる場合、観測されたデータに対してモデルが高い確率密度を与えるパラメータを選ぶ\\
\end{block}

\end{frame}
%=======================================================
\begin{frame}{Forward KL divergence}
\begin{block}{定義:KL divergence}
    \begin{align*}
        D_{KL}(p_{data}\mid\mid p_{\phi})
        &:= \int p_{data}(\boldsymbol{x})\log\frac{p_{data}(\boldsymbol{x})}{p_{\phi}(\boldsymbol{x})}\,dx \\
        &= \mathbb{E}_{\boldsymbol{x}\sim p_{data}}
            [\log p_{data}(\boldsymbol{x})-\log p_{\phi}(\boldsymbol{x})]
        \end{align*}
        二つの確率分布の隔たりを測る指標
\end{block}

\begin{itemize}
    \item 非対称\footnote{||は二つの確率分布を区別し、慣用的に非対称を表すこともある}:$D_{KL}(p_{data}\mid p_{\phi} )\neq D_{KL}(p_{\phi}\mid\mid p_{data})$
    \item 非負性:$D_{KL}(p_{data}\mid\mid p_{\phi})\geq 0$
    \item 一致時に0:$D_{KL}(p_{data}\mid\mid p_{\phi})=0$⇔$p_{data}=p_{\phi}$
\end{itemize}

\end{frame}
%=====================================================
\begin{frame}{さらに分解すると}
    \begin{block}{}
        一見すると、データの確率密度 \(p_{\mathrm{data}}(x)\) が未知であるため、\\
        順方向のKLダイバージェンスは利用できなそうだが
        \begin{align*}
            D_{\mathrm{KL}}(p_{\mathrm{data}}\mid\mid p_{\phi})
            &= \mathbb{E}_{x\sim p_{\mathrm{data}}}
                \left[
                    \log\frac{p_{\mathrm{data}}(x)}{p_{\phi}(x)}
                \right] \\
            &= -\mathbb{E}_{x\sim p_{\mathrm{data}}}
                \bigl[\log p_{\phi}(x)\bigr]
                - H(p_{\mathrm{data}}).
        \end{align*}
        ここで
        \[
            H(p_{\mathrm{data}})
            :=
            -\mathbb{E}_{x\sim p_{\mathrm{data}}}
            \bigl[\log p_{\mathrm{data}}(x)\bigr].
        \]
        \small はデータ分布のエントロピーと呼ばれ、パラメーター$\phi$に依存しない定数
    \end{block}
\end{frame}
